\begin{textsample}{POS Dim 1 – 2022\_c6 – Score -2.00 – 2022\_c6/enade\_2022\_tec\_des\_mod\_5d9084e355.txt}  \label{ex:f1_pos_266}
CAPÍTULO 7 ANÁLISE TÉCNICA DA PROVA Este capítulo tem por objetivo apresentar o desempenho dos estudantes concluintes de Tecnologia em Design de Moda no Enade/2022.
Para isso, foram calculadas as estatísticas básicas da prova em seu todo ( Subseção 7.1.1 ), bem como as estatísticas relacionadas aos componentes de formação geral ( Subseção 7.1.2 ) e de conhecimento específico da área ( Subseção 7.1.3 ).
Dadas as suas características, foram analisadas, em separado, as questões objetivas ( Seção 7.2 ) e as questões discursivas ( Seção 7.3 ).
Tomando -se como base as duas questões discursivas do Componente de Formação Geral, nas Subseções 7.3.1.6 a 7.3.1.8, são apresentados os resultados e os comentários sobre a correção das respostas em relação à Língua Portuguesa.
Para as questões objetivas, foram disponibilizados os índices de Facilidade e de Discriminação Ponto-Bisserial, também em separado, para os Componentes de Formação Geral ( Subseção 7.2.1 ) e de Conhecimento Específico ( Subseção 7.2.2 ).
De cada componente, uma das questões foi escolhida para exemplificar a análise gráfica, relacionando -se as alternativas escolhidas pelos estudantes ( inclusive o gabarito ) com o número de acertos no componente.
No Anexo IV, apresenta -se a íntegra da análise gráfica para todas as questões objetivas.
Para cada uma das questões discursivas, os conteúdos dos tipos mais comuns de respostas dos estudantes são apresentados e comparados com o padrão de resposta esperado ( ver Anexo IX ).
Nas tabelas, constam as seguintes estatísticas das notas¹¹ : média do desempenho na prova, erro-padrão da média, desvio-padrão, nota mínima, mediana e nota máxima para cada um de seus componentes.
Tais estatísticas contemplam o total de estudantes concluintes da área de Tecnologia em Design de Moda inscritos e presentes à prova do Enade/2022, tendo em vista agregações, ou por grandes regiões e o país como um todo, ou por categoria administrativa, organização acadêmica da IES e modalidade de oferta dos cursos.
Em relação aos gráficos de distribuição de notas, o intervalo considerado foi de 10 unidades, aberto à esquerda e fechado à direita, com exceção do primeiro intervalo, [ 0 ; 10 ], fechado em ambos os ¹¹ Uma definição dessas estatísticas pode ser encontrada no Glossário inserido no final desse relatório. extremos.
Para os gráficos de distribuição das notas das questões discursivas, foram consideradas mais duas categorias : questão em branco12 e nota zero. 7.1 ESTATÍSTICAS BÁSICAS DA PROVA Esta seção apresenta estatísticas selecionadas e histogramas da nota geral ( 7.1.1 ) e de cada componente : formação geral ( 7.1.2 ) e conhecimento específico ( 7.1.3 ).
São também apresentadas estatísticas selecionadas de subpopulações, caracterizadas por grande região, categoria administrativa, organização acadêmica e modalidade de oferta. 7.1.1 ESTATÍSTICAS BÁSICAS GERAIS Na Tabela 7.1, são apresentadas as estatísticas básicas da prova, por grande região, dos estudantes concluintes de Tecnologia em Design de Moda.
Como se pode verificar, a partir dos dados apresentados na Tabela 7.1, A Média das notas da prova, como um todo, foi 60,0.
Os estudantes da região Norte obtiveram a Média mais baixa ( 55,2 ) e os da região Nordeste alcançaram a Média mais alta ( 62,5 ).
As demais Médias foram : 61,3, na região Centro-Oeste ; 59,9, na região Sudeste ; e 59,4, na região Sul.
O Desvio-padrão para o Brasil, como um todo, foi de 14,8, sendo o maior valor encontrado na região Sul ( 15,0 ), e o menor, na região Nordeste ( 4,0 ), indicando uma menor dispersão das notas desta última região.
A região que obteve a maior nota “ Máxima ” foi a Sul ( 92,0 ), ao passo que a região que atingiu a menor nota “ Máxima ” foi a Norte ( 59,8 ).
A “ Mediana ” do Brasil, como um todo, foi de 61,7, sendo a maior “ Mediana ” obtida na região Nordeste ( 65,2 ), e a menor obtida na região Norte ( 54,9 ).
A única região com nota “ Mínima ” igual a zero foi a região Sul, as outras regiões obtiveram 13,2 no Centro-Oeste, 17,5 na região Nordeste, 51,3 na região Norte, e 15,1 na região Sudeste. nota “ Mínima ” igual a zero foi a região Sul, as outras regiões obtiveram 13,2 no Centro-Oeste, 17,5 na região Nordeste, 51,3 na região Norte, e 15,1 na região Sudeste.
O comportamento das notas dos estudantes de todo o Brasil pode ser observado no Gráfico 7.1, que apresenta um histograma com a distribuição das mesmas : uma distribuição unimodal com a moda no intervalo ( 50 ; 60 ].
Como mostra o Gráfico 7.1, a moda da distribuição ocorre no intervalo ( 60 ; 70 ], com frequência de aproximadamente 30 \% dos participantes.
A distribuição é crescente até o intervalo modal e diminui conforme o aumento das notas.
Na Tabela 7.2, são apresentadas informações referentes à “ Média ” da nota final, desagregadas por categoria administrativa, organização acadêmica e modalidade de oferta dos cursos.
Os dados da Tabela 7.2 mostram que, em relação à categoria administrativa, os estudantes das IES públicas obtiveram Média mais alta ( 61,5 ) que a Média das IES privadas ( 59,8 ).
O Desvio-padrão foi 14,8 tanto para as IES públicas quanto para as privadas. 7.1.2 ESTATÍSTICAS BÁSICAS NO COMPONENTE DE FORMAÇÃO GERAL Na Tabela 7.3, são apresentadas as estatísticas básicas em relação ao componente da prova que avalia a formação geral dos estudantes concluintes.
A nota de formação geral é obtida levando -se em conta os acertos nas oito questões objetivas e a nota média obtida nas questões discursivas 1 e 2, sendo a nota de cada questão discursiva obtida levando -se em conta 80 \% da nota relativa ao conteúdo e 20 \% da nota relativa ao desempenho linguístico.
Conforme os dados da Tabela 7.3, os estudantes de todo o Brasil obtiveram desempenho médio de 49,5.
Quanto à variabilidade, o Desvio-padrão das notas dos estudantes do Brasil foi 18,2.
A maior Média foi obtida na região Nordeste ( 54,0 ) e a menor, na região Norte ( 41,2 ).
As demais Médias foram : 51,0, na região Centro-Oeste, 50,9, na região Sudeste e 48,0, na região Sul.
Já o maior Desvio-padrão foi obtido na região Nordeste ( 19,0 ), e o menor, na região Norte ( 17,0 ).
Os demais Desvios foram : 17,5 na região Sudeste, 18,1 na região Sul, e 18,0 na região Centro-Oeste.
O comportamento das notas do Componente de Formação Geral pode ser observado no Gráfico 7.2, que apresenta um histograma com a distribuição dessas : uma distribuição unimodal com a moda no intervalo ( 50 ; 60 ].
Nota -se que, no Gráfico 7.2, as notas apresentam dispersão um pouco maior do que no Gráfico 7.1, anteriormente mostrado ( distribuição das notas da prova ), confirmada pela comparação dos desvios padrões : 14,8 para a nota da prova, como um todo, e 18,8 para o Componente de Formação Geral.
Na Tabela 7.4, são apresentadas as informações referentes ao desempenho dos concluintes do Componente de Formação Geral, em diferentes agregações : categoria administrativa, organização acadêmica e modalidade de oferta.
A partir dos dados da Tabela 7.4, observa -se que a maior Média foi obtida por estudantes de IES públicas ( 50,3 ) e a menor, pelos estudantes de IES privadas ( 49,4 ).
Já o maior Desvio-padrão foi obtido nas instituições privadas ( 18,4 ). 7.1.3 ESTATÍSTICAS BÁSICAS DO COMPONENTE DE CONHECIMENTO ESPECÍFICO Na Tabela 7.5, são apresentadas as estatísticas básicas referentes ao Componente Específico da área de Tecnologia em Design de Moda.
A nota do componente de conhecimentos específicos leva em conta as notas da parte objetiva e a média das notas obtidas nas três questões discursivas parte discursiva da área.
De acordo com a Tabela 7.5, a Média do desempenho dos estudantes do Brasil, como um todo, foi de 63,5.
A maior Média foi obtida na região Nordeste ( 65,3 ), e a menor, na região Norte ( 59,9 ).
As demais Médias foram : na região Centro-Oeste, 64,7 ; na região Sudeste 63,0 ; e na região Sul, 63,3.
Quanto à variabilidade das notas, o Desvio-padrão do Brasil, como um todo, foi de 15,9, sendo o maior Desvio- padrão observado na região Sul ( 16,1 ) e o menor, na região Norte ( 4,4 ).
Os demais desvios foram : 15,3 na região Centro-Oeste, 16,0 na região Sudeste e 15,7 na região Nordeste.
A “ Mediana ” das notas dos estudantes de todo o Brasil foi 65,5.
A maior “ Mediana ” ocorreu na região Nordeste ( 67,5 ), e a menor, na região Norte ( 59,4 ).
As demais “ Medianas ” foram : 66,1 na região Centro-Oeste, 65,8 na região Sudeste e 64,8 na região Sul.
A nota “ Máxima ” do Brasil, como um todo, foi de 99,0, sendo obtida por, pelo menos, um estudante na região Sudeste.
As demais notas “ Máximas ” foram : 65,6 na região Norte, 96,9 na região Nordeste, 96,4 na região Sul e 93,7 na região Centro-Oeste.
A nota “ Mínima ” foi zero apenas na região Sul ; nas demais, 12,6 na região Centro-Oeste, 11,4 na região Nordeste, 54,0 na região Norte e 12,6 na região Sudeste.
O comportamento das notas do Componente Específico pode ser observado no Gráfico 7.3, que apresenta um histograma com a distribuição dessas : uma distribuição unimodal com a moda no intervalo ( 60 ; 70 ].
Na Tabela 7.6, é apresentada uma comparação dos resultados em relação a categoria administrativa, a organização acadêmica e a modalidade de oferta, agora levando -se em conta o desempenho de estudantes do Componente Específico da prova.
Quanto à categoria administrativa, observa -se a maior Média foi obtida por estudantes das IES públicas ( 65,2 ) e a menor, pelos estudantes de IES privadas ( 63,2 ). 7.2 ANÁLISE DAS QUESTÕES OBJETIVAS Esta seção apresenta estatísticas selecionadas e histogramas das questões objetivas de formação geral ( 7.2.1 ) e de conhecimento específico ( 7.2.2 ).
São, também, apresentadas e comparadas as médias das subpopulações caracterizadas por grande região.13 7.2.1 COMPONENTE DE FORMAÇÃO GERAL Na Tabela 7.7, são apresentadas as estatísticas básicas relativas às oito questões objetivas do componente da prova que abrange a formação geral dos estudantes. 13 As fórmulas utilizadas para o coeficiente de correlação ponto-bisserial e para o índice de facilidade estão apresentadas no Capítulo 2.
De acordo com os dados da Tabela 7.7, a “ Média ” do Brasil foi 51,3.
A menor “ Média ” foi encontrada na região Norte ( 39,6 ), e a maior, na região Nordeste ( 56,0 ).
As demais “ Médias ” foram : 53,4 na região Sudeste, 49,0 na região Sul e 55,7 na região Centro-Oeste.
O “ Desvio-padrão ” do Brasil foi 21,3, sendo o maior “ Desvio-padrão ” encontrado na região Nordeste ( 22,9 ), e o menor “ Desvio-padrão ” encontrado na região Norte ( 16,6 ).
As “ Medianas ” do Brasil, como um todo, e de três regiões ( Centro-Oeste, Sudeste e Sul ) foram iguais a 50,0.
Na região Nordeste, a “ Mediana ” foi maior, igual a 62,5 e na região Norte, foi menor, 37,5.
A nota “ Máxima ” ( 100,0 ) foi alcançada em quase todas as grandes regiões, com exceção do Norte ( 62,5 ).
As notas “ Mínimas ” ( 0,0 ) foram iguais para todas as regiões, menos para o Norte ( 25,0 ).
Na Tabela 7.8, exposta mais adiante, são apresentados o Índice de Facilidade e o Índice de Discriminação ( Ponto-Bisserial ) para cada uma das questões objetivas do Componente de Formação Geral.
De acordo com os dados da Tabela 7.8, as questões objetivas do Componente de Formação Geral, segundo o Índice de Facilidade, foram assim avaliadas : apenas a questão 3 teve o Índice de Facilidade classificado como “ Muito fácil ” por ter um índice de acerto de 0,86.
Três questões foram consideradas com índice de dificuldade “ Médio ” por terem índice de acerto entre 0,48 e 0,52.
Três questões foram classificadas como difíceis, por terem índice de acerto entre 0,39 e 0,49.
Já o Índice de Discriminação ( Ponto-Bisserial ) variou entre “ Bom ” e “ Muito bom ”.
Na Tabela 7.9, é apresentada a distribuição das questões, levando -se em conta, simultaneamente, a classificação dos dois índices.
Os dados da Tabela 7.9 atestam que seis questões que tiveram Índice de Discriminação “ Muito bom ” figuraram em três níveis de dificuldade : “ Fácil ”, “ Médio ” e “ Difícil ”.
Três classificadas na categoria “ Médio ” do Índice de Facilidade, duas na categoria “ Difícil ” e uma na categoria “ Fácil ”.
As questões que tiveram Índice de Discriminação “ Bom ” figuraram no nível de dificuldade “ Difícil ” e “ Muito fácil ”.
No Gráfico 7.4, para exemplificar, analisa -se o comportamento da questão de número 6 de formação geral.
Trata -se de uma questão considerada com nível “ Médio ” de dificuldade, em relação à facilidade, e que obteve o Índice de Discriminação 0,51 ( “ Muito bom ” ). consideradas com índice de dificuldade “ Médio ” por terem índice de acerto entre 0,48 e 0,52.
Três questões foram classificadas como difíceis, por terem índice de acerto entre 0,39 e 0,49.
Já o Índice de Discriminação ( Ponto-Bisserial ) variou entre “ Bom ” e “ Muito bom ”.
Na Tabela 7.9, é apresentada a distribuição das questões, levando -se em conta, simultaneamente, a classificação dos dois índices.
Os dados da Tabela 7.9 atestam que seis questões que tiveram Índice de Discriminação “ Muito bom ” figuraram em três níveis de dificuldade : “ Fácil ”, “ Médio ” e “ Difícil ”.
Três classificadas na categoria “ Médio ” do Índice de Facilidade, duas na categoria “ Difícil ” e uma na categoria “ Fácil ”.
As questões que tiveram Índice de Discriminação “ Bom ” figuraram no nível de dificuldade “ Difícil ” e “ Muito fácil ”.
No Gráfico 7.4, para exemplificar, analisa -se o comportamento da questão de número 6 de formação geral.
Trata -se de uma questão considerada com nível “ Médio ” de dificuldade, em relação à facilidade, e que obteve o Índice de Discriminação 0,51 ( “ Muito bom ” ).
No Gráfico 7.4, cada uma das cinco curvas representa o percentual de respostas em determinada alternativa da questão, em função do número de acertos dos estudantes nessa parte da prova ( formação geral/múltipla escolha ), antes de possíveis eliminações pelo critério Ponto-Bisserial.
Os gráficos relativos às demais questões de formação geral constam do Anexo IV. 7.2.2 COMPONENTE ESPECÍFICO Na Tabela 7.10, são apresentadas as estatísticas básicas em relação às questões objetivas do Componente Específico da prova de Tecnologia em Design de Moda por grande região.
Como se verifica a partir dos dados da Tabela 7.10, a “ Média ” do Brasil do Componente Específico foi 64,4.
A menor “ Média ” foi observada na região Norte ( 61,8 ) e a maior, na região Nordeste ( 66,1 ).
O Desvio-padrão de todo o Brasil foi 16,1, sendo o menor Desvio-padrão encontrado na região Norte ( 7,6 ) e o maior, na região Sudeste ( 16,7 ).
A “ Mediana ” de todo o Brasil foi 66,7, o mesmo valor da “ Mediana ” encontrada nas regiões Centro-Oeste, Nordeste, Sudeste e Sul.
Na região Norte, a “ Mediana ” foi 61,2.
A nota “ Máxima ” da prova foi 100,00, obtida nas questões objetivas do Componente Específico, por, pelo menos, um estudante das regiões Sudeste, Sul e Nordeste.
Na região Centro-Oeste, o valor foi 92,6 e na região Norte, esse valor foi 70,4.
A nota “ Mínima ” foi 0 no Brasil e no Sul, 14,8 no Centro-Oeste, 7,4 no Nordeste, 51,9 no Norte e 11,1 no Sudeste.
Na Tabela 7.11, são apresentados os Índices de Facilidade e de Discriminação ( Ponto-Bisserial ) das questões objetivas do Componente Específico, para os estudantes de Tecnologia em Design de Moda.
A partir do Índice de Facilidade obtido, apresentado na Tabela 7.11, evidenciamos apenas duas questões consideradas como “ muito fácil ”, que são as questões 12 e 16, consideradas “ muito bom ” e “ médio ”, respectivamente ; seguidas por 14 questões consideradas fáceis, das quais duas questões tiveram o índice de discriminação “ muito bom ” ( 15 e 29 ) e uma questão considerada “ médio ” ( 35 ).
Em menor quantidade, há questões consideradas “ difícil ”, apenas uma ao todo ( 33 ) e com índice de discriminação “ médio ”.
Na Tabela 7.12, é apresentada a distribuição das questões levando -se em conta, simultaneamente, a classificação dos dois índices.
Como revelam os dados da Tabela 7.12, cinco questões que tiveram Índice de Discriminação “ muito bom ” figuraram em três níveis de dificuldade, “ muito fácil ”, “ fácil ” e “ médio ” : duas classificadas na categoria “ médio ” do Índice de Facilidade, duas na categoria “ fácil ” e uma na categoria “ muito fácil ”.
As questões que tiveram Índice de Discriminação bom figuraram no nível de facilidade “ médio ” ( 6 ) e “ fácil ” ( 11 ).
As questões que tiveram Índice de Discriminação “ médio ” figuraram no nível de facilidade “ muito fácil ” ( 1 ), “ fácil ” ( 1 ), “ médio ” ( 1 ) e “ difícil ” ( 1 ).
Houve apenas uma questão com dificuldade “ difícil ”.
No Gráfico 7.5, cada uma das cinco curvas representa o percentual de respostas em determinada alternativa da questão, em função do número de acertos dos estudantes nessa parte da prova ( Componente Específico ), antes de possíveis eliminações pelo critério Ponto-Bisserial.
Os gráficos relativos às demais questões do Componente Específico constam no Anexo IV. 7.3 ANÁLISE DAS QUESTÕES DISCURSIVAS Esta seção apresenta estatísticas selecionadas e histogramas das questões discursivas de formação geral ( 7.3.1 ) e de Componente Específico ( 7.3.2 ).
Expõe e compara, também, as médias de subpopulações, caracterizadas por grande região.
Cumpre notar que nem todas as questões passam pelo crivo dos corretores.
As respostas depois de digitalizadas passam por um pré-processamento de identificação automático.
Algumas destas questões são definidas como um “ branco automático ” ( que não são enviadas para os professores para correção, a não ser que na segunda fase, a de inspeção visual, seja constatada uma resposta ).
O critério usado neste pré-processamento, é verificar para cada questão, a partir da quantidade de pixels, o que provavelmente, é uma resposta em branco.
A partir da separação destes documentos que foram considerados em branco, eles são enviados para uma verificação visual, por uma equipe de colaboradores, devidamente capacitados.
Deste modo, quaisquer erros, deste pré-processamento, são encontrados, e as questões preenchidas são enviadas para a correção dos professores.
Eventualmente, algumas respostas em branco não são detectadas nesse pré-processamento automático.
A presença de pequenos riscos, ou manchas, pode fazer o programa não identificar uma questão como um possível branco.
Estas são enviadas para a correção, mas são classificadas como “ branco ” pelos professores, quando da correção. 7.3.1 COMPONENTE DE FORMAÇÃO GERAL As análises dos resultados de desempenho dos estudantes de Tecnologia em Design de Moda, nas duas questões discursivas relativas à formação geral, encontram -se na Tabela 7.13 e no Gráfico 7.6, cujos dados são analisados de forma sequenciada.
No caso da Tabela 7.13, analisa -se a média das notas de cada questão discursiva de formação geral que leva em conta as correções realizadas por duas bancas : uma que analisa e pontua aspectos relativos ao conteúdo das respostas, seguindo o padrão de respostas divulgado pelo Inep ( ver Anexo IX ) ; outra que avalia o desempenho linguístico expresso pela redação das respostas, seguindo um padrão de respostas específico da área de Língua Portuguesa, detalhado na seção 7.3.1.7 deste relatório.
A nota de cada questão é composta de 80 \% da nota da correção de conteúdo e 20 \% da nota de desempenho linguístico.
Na Tabela 7.13, observa -se que a nota “ Média ” nesse conjunto de questões foi inferior à obtida nas objetivas.
Os estudantes de todo o Brasil obtiveram, em formação geral, média 46,9 nas questões discursivas.
Pode -se notar, também, que o “ Desvio-padrão ” nesse conjunto de questões foi de 24,7 nas A maior “ Média ” foi obtida na região Nordeste ( 51,0 ) e a menor, na região Norte ( 43,7 ).
A “ Mediana ” de todo o Brasil, neste componente, foi 52,3.
Nas regiões Sul, Sudeste e Centro-Oeste, a “ Mediana ” foi menor ( 52,0, 51,0 e 50,2 respectivamente ), e, nas demais regiões, foi maior que o valor para o Brasil : Norte ( 53,0 ) e Nordeste ( 56,5 ).
A nota “ Máxima ” ( 100,0 ) foi obtida na região Sul, tendo sido 70,5 a “ Máxima ” na região Norte ; 94,5, na região Nordeste ; 93, na região Sudeste ; e 94,5, na região Centro-Oeste.
A nota “ Mínima ” ( 0,0 ) foi a mesma em todas as regiões do Brasil.
No Gráfico 7.6, está representada a distribuição das notas nas questões discursivas do Componente de Formação Geral.
Relatório Síntese de Área – ENADE/2022 Como mostra o Gráfico 7.6, o comportamento das notas das questões discursivas de formação geral apresenta um histograma com a distribuição das mesmas : uma distribuição unimodal com a moda no intervalo ( 60 ; 70 ].
Na sequência, os resultados verificados para cada uma das questões discursivas de formação geral serão apresentados, estabelecendo -se relações com os conteúdos abordados em cada uma delas.
Os comentários das bancas de docentes corretores, a respeito do observado na correção das respostas dos estudantes, suas impressões e conclusões serão apresentados junto à análise de cada questão.
Cumpre esclarecer que, tendo em vista que as questões discursivas de formação geral são padronizadas, ou seja, constam de todas as provas, os comentários da Banca são os mesmos para todas as carreiras acadêmicas, sendo direcionados a todos os estudantes que participaram do Enade/2022.
A seguir, serão analisados o desempenho linguístico dos estudantes da área de Tecnologia em Design de Moda nas duas questões discursivas de formação geral do Enade/2022, apresentando -se os resultados obtidos e os comentários da banca de correção para cada questão. 7.3.1.1 ANÁLISE DE CONTEÚDO DA QUESTÃO DISCURSIVA 1 DO COMPONENTE DE FORMAÇÃO GERAL Os dados de Tecnologia em Design de Moda obtidos a partir das respostas à Questão 1 do Componente de Formação Geral, encontram -se na Tabela 7.14 e no Gráfico 7.7, exposto após a apresentação dos dados da referida tabela.
Verifica -se, a partir dos dados da Tabela 7.14, que na Questão Discursiva 1 do Componente de Formação Geral, os estudantes de todo o Brasil tiveram “ Média ” de 43,8.
A maior “ Média ” para a Questão 1 foi obtida na região Norte ( 50,1 ), e a menor, na região Centro-Oeste ( 42,3 ).
Quanto à variabilidade das notas, o “ Desvio-padrão ” de todo o Brasil foi 26,5.
O menor “ Desvio-padrão ” foi obtido na região Sudeste ( 26,0 ), e o maior “ Desvio-padrão ” foi obtido pela região Nordeste ( 29,3 ).
A “ Mediana ” do Brasil, como um todo, foi de 46,5, sendo a maior obtida no Norte ( 58,3 ), e a menor na região Centro-Oeste ( 43,5 ).
A nota “ Máxima ” da Questão Discursiva 1 foi 100,0, obtida por, pelo menos, um estudante nas regiões Sudeste e Sul.
Na região Centro-Oeste, a nota “ Máxima ” foi 95,5, e a menor foi a da região Nordeste ( 75,5 ).
Já a nota “ Mínima ” para todas as regiões foi zero.
No Gráfico 7.7, apresenta -se a distribuição das notas na Questão Discursiva 1 do Componente de Formação Geral.
Como mostra o Gráfico 7.7, o comportamento das notas das Questões 1 de formação geral apresenta um histograma com a distribuição dessas : uma distribuição unimodal com a moda nos intervalos [ 0 ; 10 e 50 ; 60 ]. 7.3.1.2 COMENTÁRIOS SOBRE A CORREÇÃO DE CONTEÚDO DAS RESPOSTAS À QUESTÃO DISCURSIVA 1 A Questão 1 abordou um tema atual, com dois itens para os estudantes responderem : a ) “ discuta por que os investimentos em educação, ciência, tecnologia e inovação são estratégicos e contribuem para o desenvolvimento de um país ” ; b ) “ explique como o fomento público ao desenvolvimento científico pode atender à justiça social e à equidade, em contextos como o da pandemia de Covid-19 ”.
O tema contempla como os investimentos em educação, ciência, tecnologia e inovação ( ECTI ) cooperam para um país se desenvolver e como o fomento público, vinculado ao desenvolvimento científico, pode atender à justiça social, à equidade e à população em momentos de pandemia.
Esse tema é atual, pois o mundo vivenciou recentemente uma pandemia que, somente no Brasil, deixou mais de 700 mil mortos, sendo exaustivamente comentada na televisão e nos meios digitais.
Logo, almeja -se que os respondentes justifiquem a necessidade de investimentos em ECTI com foco em questões como : investimentos públicos em P\&D por meio de Universidades públicas ; pesquisas realizadas em outros países que possibilitaram o acesso mais célere às vacinas e serviram de modelo para o que foi desenvolvido no Brasil ; pesquisadores de universidades, tanto brasileiras quanto estrangeiras como protagonistas para orientar e informar a respeito da Covid-19 por meio dos canais de comunicação ; desenvolvimento de equipamentos de segurança, como máscaras de nariz ou de rosto, respiradores, que foram implementados e viabilizados por pesquisa e extensão em Universidades e escolas técnicas.
Esperava -se que os candidatos soubessem o básico a respeito do tema, pois o assunto é bastante atual.
Quanto ao item B, esperava -se que os candidatos abordassem questões como : o papel das Universidades e dos centros de pesquisa para a justiça e a equidade ; o investimento público em pesquisa, tecnologia e educação para garantir acesso igualitário da população aos recursos de saúde, podendo fazer um contraponto com os investimentos privados ; a política de distribuição de vacinas gratuitas por meio do Sistema Único de Saúde, que possibilitou o acesso às vacinas no Brasil ; e as questões econômicas mundiais que determinaram uma distribuição desigual de recursos de saúde em escala global, com os países mais ricos retendo a maior parte desses recursos.
Atualmente, é comprovado que os países dominantes são aqueles que se desenvolvem e dominam as ferramentas tecnológicas.
As pesquisas mostram que os países que, mesmo em crise, investiram em ciência, tecnologia e inovação são os que conseguiram se reerguer mais rapidamente e sobressair -se aos demais países em crises pandêmicas.
Portanto, a questão foi elaborada de maneira oportuna e adequada ao universo do estudante universitário brasileiro.
No item A, o estudante poderia focar em diversos argumentos, entre eles : como ocorreu o investimento público em pesquisa e desenvolvimento por meio de Universidades públicas e institutos de pesquisa ; como as pesquisas realizadas em outros países possibilitaram acesso mais rápido às vacinas e serviram de modelo para o que foi desenvolvido no Brasil ; como o contingenciamento do Fundo Nacional de Desenvolvimento Científico e Tecnológico ( FNDCT ) prejudicou e ainda prejudica o Brasil, entre diversos outros assuntos.
No item B, era esperado que o estudante explicasse como o fomento público ao desenvolvimento científico pode atender à justiça social e à equidade, em contextos como o da pandemia de Covid-19.
Além disso, deveria abordar a importância das Universidades com a justiça social ; como os países que investem nessa área têm grandes vantagens em relação a outros países que não investem ; a distribuição de vacina atende à justiça social por possibilitar que toda a população seja vacinada, independentemente de ser rico ou pobre, por exemplo.
O espaço disponível para a resposta dos dois itens foi adequado e a maioria das respostas ocupou parte das linhas, não utilizando todo o espaço.
Algumas respostas não apresentaram a separação dos itens A e B.
A maioria dos estudantes respondeu de modo semelhante ao que constava do padrão de resposta.
Foi recorrente a exposição do problema sem desenvolvê -lo com clareza, apenas de forma superficial.
A abordagem do candidato sobre o tema tornou -se, em alguns casos, muito genérica.
As respostas consideradas boas atenderam ao que a banca esperava.
Abordaram a importância de investir em CT\&I para desenvolver o país, desdobrando os argumentos e citando como isso poderia ocorrer.
No geral, as notas estavam abaixo do esperado.
As respostas evidenciaram que a maioria dos estudantes possui um conhecimento sobre os dois itens apresentados.
A Questão 1 visou avaliar o desenvolvimento da argumentação de um tema relevante para a sociedade brasileira vigente.
Além disso, houve problemas estruturais nos textos, pois a grande maioria dos estudantes tem uma noção do assunto, mas muitos possuem dificuldade de construir um texto com coesão e coerência, ou seja, não conseguem concatenar os argumentos para defender um ponto de vista.
No geral, em relação à Questão 1, identificou -se que a principal deficiência na formação dos estudantes está relacionada ao aprofundamento do conteúdo e à capacidade de expressar, com clareza e coesão, sua opinião a respeito de determinado assunto. 7.3.1.3 ANÁLISE DE CONTEÚDO DA QUESTÃO DISCURSIVA 2 DO COMPONENTE DE FORMAÇÃO GERAL A Tabela 7.15 apresenta dados estatísticos a respeito das notas de conteúdo da Questão Discursiva 2 do Componente de Formação Geral, organizados por grande região de funcionamento dos cursos, no país.
A Tabela 7.15 mostra que o desempenho médio dos estudantes, na Questão Discursiva 2 ( “ Média ” 50,0 ), foi acima do obtido na Questão Discursiva 1 ( “ Média ” 43,8 ).
Na região Norte, a “ Média ” foi menor ( 37,3 ), e a maior “ Média ” foi na região Nordeste ( 54,3 ).
Quanto à variabilidade das notas, o desvio-padrão de todo o Brasil foi 28,5, superior ao obtido na Questão Discursiva 1 ( 26,5 ).
O maior desvio nessa questão foi obtido na região Norte ( 30,6 ), enquanto o menor foi obtido na região Sudeste ( 27,4 ).
A “ Mediana ” de todo o Brasil foi 56,0, a mesma da região Sul.
A menor “ Mediana ” obtida foi na região Norte, com 47,8.
A nota “ Máxima ” da Questão Discursiva 2 foi 100,0, obtida por, pelo menos, um estudante nas regiões Sudeste e Sul.
Na região Norte, a nota “ Máxima ” foi 65,5.
Já a nota “ Mínima ” para todas as regiões foi zero.
No Gráfico 7.8, mostra -se a distribuição das notas na Questão Discursiva 2 do Componente de Formação Geral.
Como mostra o Gráfico 7.8, o comportamento das notas das Questões 2 de formação geral apresenta um histograma com a distribuição das mesmas : uma distribuição unimodal com a moda no intervalo ( 70 ; 80 ]. 7.3.1.4 COMENTÁRIOS SOBRE A CORREÇÃO DE CONTEÚDO DAS RESPOSTAS À QUESTÃO DISCURSIVA 2 A Questão Discursiva 2 abordou um tema importante, com o objetivo de analisar conhecimentos gerais do aluno.
O tema contemplou o conhecimento geral do patriarcado estrutural ( dominação masculina ), o qual é recorrente nos noticiários, e um infortúnio na vida dos brasileiros, pois trata -se de uma questão de violência contra as mulheres.
É recorrente assistir ou ler notícias sobre o domínio masculino na sociedade brasileira.
Questões como a participação da mulher na política ; a paridade de salários entre profissionais de sexos diferentes ; a desigualdade social e o tratamento hostil que é dado às mulheres, em diversas situações, tanto profissionais quanto familiares, devem ser de conhecimento geral do estudante.
Não se objetiva tratar questões de ordem específica para a solução deste problema, mas, apenas, discorrer e argumentar sobre ele.
Portanto, a questão foi elaborada de maneira oportuna e adequada para o universo do estudante universitário brasileiro.
A Questão está dividida em dois itens : A e B.
No item A, esperava -se que o estudante abordasse o patriarcado – dominação masculina que gera a divisão sexual do trabalho, a desigualdade econômica e política, entre homens e mulheres, a sub-representação das mulheres na política e em cargos de poder.
Esperava -se que o estudante ressaltasse, ainda, questões culturais, como a educação sexista e a mídia, que coisificam a imagem das mulheres e naturalizam a desigualdade entre os sexos, de forma a considerar os homens como dominantes, enquanto as mulheres são historicamente tidas como submissas e incapazes.
No item B, era esperado que o estudante apresentasse a importância do papel do Estado na implementação de políticas públicas voltadas à proteção da mulher, tais como na área da saúde, na educação não sexista, na assistência social, na segurança pública – delegacias femininas, defensoria pública especializada, casa de abrigo, juizado especial etc., bem como ONGs que atuam no combate à violência contra a mulher e ao feminicídio.
O espaço disponível para a resposta dos dois itens foi adequado e grande parte das respostas ocupou parte das linhas, não contemplando o espaço completo.
Com relação à linguagem utilizada pelos estudantes, as respostas estavam dentro do esperado.
Alguns estudantes utilizavam termos do jargão do cotidiano.
No item B, o estudante deveria propor duas ações do Estado para o enfrentamento da violência contra as mulheres.
Os alunos deveriam apresentar propostas e desenvolvê -las.
Muitos estudantes responderam apenas por meio da apresentação de ações do Estado, mas sem desenvolvê -las ou argumentá -las.
Em geral, os alunos apenas expuseram ações – algumas vezes, uma, duas e, até mesmo, três ações –, mas sem fazer relações entre elas ou sem desenvolvê -las.
As respostas dos estudantes demonstraram que a maioria apresenta um conhecimento geral a respeito do patriarcado estrutural e da trágica desigualdade que esta impõe dos homens sobre as mulheres, bem como de possíveis ações do Estado para a proteção delas.
A Questão visou avaliar o desenvolvimento da argumentação de um tema relevante para a sociedade brasileira vigente. salários entre profissionais de sexos diferentes ; a desigualdade social e o tratamento hostil que é dado às mulheres, em diversas situações, tanto profissionais quanto familiares, devem ser de conhecimento geral do estudante.
Não se objetiva tratar questões de ordem específica para a solução deste problema, mas, apenas, discorrer e argumentar sobre ele.
Portanto, a questão foi elaborada de maneira oportuna e adequada para o universo do estudante universitário brasileiro.
A Questão está dividida em dois itens : A e B.
No item A, esperava -se que o estudante abordasse o patriarcado – dominação masculina que gera a divisão sexual do trabalho, a desigualdade econômica e política, entre homens e mulheres, a sub-representação das mulheres na política e em cargos de poder.
Esperava -se que o estudante ressaltasse, ainda, questões culturais, como a educação sexista e a mídia, que coisificam a imagem das mulheres e naturalizam a desigualdade entre os sexos, de forma a considerar os homens como dominantes, enquanto as mulheres são historicamente tidas como submissas e incapazes.
No item B, era esperado que o estudante apresentasse a importância do papel do Estado na implementação de políticas públicas voltadas à proteção da mulher, tais como na área da saúde, na educação não sexista, na assistência social, na segurança pública – delegacias femininas, defensoria pública especializada, casa de abrigo, juizado especial etc., bem como ONGs que atuam no combate à violência contra a mulher e ao feminicídio.
O espaço disponível para a resposta dos dois itens foi adequado e grande parte das respostas ocupou parte das linhas, não contemplando o espaço completo.
Com relação à linguagem utilizada pelos estudantes, as respostas estavam dentro do esperado.
Alguns estudantes utilizavam termos do jargão do cotidiano.
No item B, o estudante deveria propor duas ações do Estado para o enfrentamento da violência contra as mulheres.
Os alunos deveriam apresentar propostas e desenvolvê -las.
Muitos estudantes responderam apenas por meio da apresentação de ações do Estado, mas sem desenvolvê -las ou argumentá -las.
Em geral, os alunos apenas expuseram ações – algumas vezes, uma, duas e, até mesmo, três ações –, mas sem fazer relações entre elas ou sem desenvolvê -las.
As respostas dos estudantes demonstraram que a maioria apresenta um conhecimento geral a respeito do patriarcado estrutural e da trágica desigualdade que esta impõe dos homens sobre as mulheres, bem como de possíveis ações do Estado para a proteção delas.
A Questão visou avaliar o desenvolvimento da argumentação de um tema relevante para a sociedade brasileira vigente. 7.3.1.5 CONSIDERAÇÕES FINAIS SOBRE A CORREÇÃO DE CONTEÚDO DAS QUESTÕES DISCURSIVAS DO COMPONENTE DE FORMAÇÃO GERAL As questões discursivas da área de formação geral abordam temas atuais e proporcionaram a avaliação de competências importantes dentre as previstas pelas suas diretrizes curriculares.
A análise das respostas revelou, principalmente, dificuldade dos estudantes em relação à seleção de argumentos sólidos para discorrer sobre o tema – que pode ser fruto de desinteresse ou de desconhecimento mais profundo do contexto.
Esses fatos parecem ter sido os principais responsáveis pelo resultado aquém do esperado em termos de avaliação quantitativa.
Mesmo os estudantes que responderam em conformidade com o padrão demonstraram um universo de argumentação muito limitado, pobre em vocabulário, com o recurso frequente de procurar nos textos introdutórios os elementos para a construção de suas frases.
Encadear e relacionar argumentos foi uma deficiência percebida.
As razões eram apresentadas de forma confusa, sem fluência e sem constituir unidade, utilizando termos genéricos e vagos. 7.3.1.6 ANÁLISE DE LÍNGUA PORTUGUESA DAS QUESTÕES DISCURSIVAS DO COMPONENTE DE FORMAÇÃO GERAL Os dados de Tecnologia em Design de Moda obtidos a partir das respostas às questões discursivas do Componente de Formação Geral, no que tange à Língua Portuguesa, são apresentados na Tabela 7.16.
Os dados da Tabela 7.16 mostram que, em relação ao quesito Língua Portuguesa, os estudantes de todo o Brasil obtiveram “ Média ” 56,4.
A maior “ Média ” em relação à Língua Portuguesa foi obtida na região Norte ( 58,8 ), e a menor, na região Centro-Oeste ( 55,5 ).
Quanto à variabilidade das notas, o “ Desvio-padrão ” de todo o Brasil foi 28,9.
O menor “ Desvio-padrão ” foi obtido nas regiões Sudeste e Nordeste ( 28,6 ), e o maior desvio-padrão foi obtido na região Centro-Oeste ( 29,6 ).
A “ Mediana ” das notas de Língua Portuguesa foi 65,0 para o Brasil, como um todo, a mesma obtida em todas as regiões, exceto a região Norte, que obteve “ Mediana ” de 66,3.
A nota “ Máxima ” para No Gráfico 7.9, mostra -se a distribuição das notas de Língua Portuguesa do Componente de Formação Geral.
Como mostra o Gráfico 7.9, o comportamento das notas das Questões 3 de formação geral apresenta um histograma com a distribuição dessas : uma distribuição unimodal com a moda no intervalo ( 60 ; 70 ]. 7.3.1.7 COMENTÁRIOS SOBRE A CORREÇÃO DE LÍNGUA PORTUGUESA DAS RESPOSTAS ÀS QUESTÕES DISCURSIVAS Os comentários feitos neste item a respeito da correção de Língua Portuguesa das respostas emitidas para as questões discursivas do Componente de Formação Geral estão agrupados em três eixos : avaliação de desempenho linguístico ; padrão de resposta utilizado na avaliação e análise das respostas dadas pelos estudantes.
De acordo com a Portaria n.º 298 do MEC/INEP, de 30 de junho de 2022, o Enade deve avaliar se o estudante desenvolveu, no processo de formação, dentre outras, competências para [ … ] V – Compreender as linguagens e suas respectivas variações ; VI – Ler, produzir e interpretar textos com clareza e com coerência ; VII – Identificar, analisar e interpretar diferentes representações verbais, não verbais, gráficas e numéricas de fenômenos diversos ou de um mesmo significado ; VIII – Formular e articular argumentos e contra-argumentos consistentes em situações sociocomunicativas ; IX – Identificar, compreender e analisar situações-problema utilizando pensamento holístico e sistêmico ao se abordar a complexidade da realidade. ( INEP, 2022d ) Essas competências são essenciais à formação de nível superior, porque dizem respeito ao domínio da linguagem, tanto no processo de leitura quanto na articulação e na organização de ideias para a produção textual.
Dessa forma, as duas questões discursivas de formação geral oferecem oportunidade para avaliar se o participante consegue desenvolver um tema, recorrendo a argumentos que explicitem sua opinião e elaborando propostas de intervenção social adequadas a um exercício consciente da cidadania.
É, portanto, necessário identificar, nas respostas discursivas de formação geral, as diferentes habilidades necessárias ao bom desempenho linguístico na situação de formalidade exigida : a organização das ideias, tendo em vista a construção de um texto coerente e estruturado ; o domínio das convenções ortográficas e das convenções morfossintáticas que constituem o arcabouço da gramática da norma-padrão da Língua Portuguesa ; e o domínio de um vocabulário amplo e diversificado, que permita uma comunicação clara e eficiente.
A Questão Discursiva 1 desenvolveu o tema de investimentos públicos em educação, ciência, tecnologia e inovação, que são estratégicos e contribuem para o desenvolvimento científico de um país, podendo abordar essa questão tanto no Brasil quanto em outros países.
O enunciado estava claro ao encaminhar a solicitação de resposta, evidenciando a exigência de elaboração de um texto que discutisse e explicasse por que os investimentos públicos em educação, ciência, tecnologia e inovação são estratégicos e contribuem para o desenvolvimento científico de um país e como o fomento público ao desenvolvimento científico pode atender à justiça social e à equidade, em contextos como o da pandemia de Covid-19.
Os três textos motivadores favoreceram a formulação e a articulação de argumentos e contra-argumentos consistentes, em defesa do ponto de vista a ser defendido pelo participante.
A Questão Discursiva 2 desenvolveu o conceito de patriarcado, tomando como base um texto de Biroli.
Com base nesse texto, o comando solicitou do participante : “ discorra sobre a relação entre o patriarcado estrutural e a desigualdade entre homens e mulheres ” e “ proponha duas ações do Estado para o enfrentamento dos diferentes tipos de violência contra a mulher ”.
O comando “ discorra ” demandava a elaboração de um texto de base dissertativa, com desenvolvimento de argumentação sobre o tema.
Em relação à segunda parte, o comando exigia a elaboração de duas propostas de ações, que deveriam conter informações sobre o agente da intervenção ( Estado ) para o enfrentamento da situação apresentada.
As duas questões, portanto, assemelhavam -se, pois solicitavam um texto dissertativo ao apresentarem comandos para “ discutir ” e “ discorrer ”.
A Questão Discursiva 2, entretanto, solicitava que o aluno apresentasse propostas de ações para o enfrentamento do problema destacado.
Esses dois estilos textuais ofereceram oportunidade de avaliar se os participantes estavam habilitados a desenvolverem textos com esses dois distintos objetivos comunicativos, o que seria esperado ao final do curso de graduação e necessário à sua vida profissional.
A análise do desempenho linguístico esperado de um estudante que está concluindo seu curso superior apoiou -se no princípio de que ele deveria utilizar conhecimentos sobre os temas e estruturar seus textos de acordo com as características do registro formal adequado à situação comunicativa – avaliação de conhecimentos ao final de um curso superior.
Essa configuração determina exigências quanto à adequação da seleção vocabular, ao desenvolvimento do conteúdo, à estruturação sintática dos períodos, à organização lógica das ideias, à utilização de procedimentos de encadeamento textual e de referenciação, à obediência às exigências morfossintáticas próprias da modalidade escrita da norma-padrão, ao respeito às regras ortográficas de acentuação gráfica e de grafia das palavras. a ) O padrão de resposta utilizado na avaliação Em relação a este eixo de análise, foram considerados, no padrão de resposta, as seguintes competências : ( 1 ) estruturação textual ; ( 2 ) respeito às convenções ortográficas da norma-padrão da Língua Portuguesa ; ( 3 ) domínio dos diferentes aspectos morfossintáticos próprios da modalidade escrita formal da norma-padrão da Língua Portuguesa ; e ( 4 ) seleção vocabular adequada à modalidade escrita formal da Língua Portuguesa, exigida pela situação comunicativa.
A seguir, são apresentados os aspectos avaliados associados a cada uma destas competências. ( 1 ) Estruturação textual condizente com o gênero textual solicitado e o modo de organização textual adequado ao gênero ( argumentativo ou expositivo ).
Esta competência envolve : - estruturação sintática condizente com o padrão da modalidade escrita formal da Língua Portuguesa, de modo a garantir a clareza e a coerência necessárias ; - distribuição do conteúdo do texto em parágrafos, de modo a garantir a sua organização temática ; - distribuição das ideias em períodos que mantenham uma estruturação sintática condizente com a modalidade formal da língua ( evitando -se frases fragmentadas, sequência justaposta de ideias sem encaixamentos sintáticos e presença de oração subordinada sem oração principal que reflitam hábitos da oralidade e comprometam a estrutura lógico-gramatical do texto ) ; - utilização de operadores discursivos que contribuam para a progressão temática do texto, estabelecendo relações lógicas entre as ideias apresentadas, tanto do ponto de vista intrafrasal, como interfrasal : preposições, conjunções, pronomes relativos, advérbios e locuções adverbiais de tempo e lugar ; - utilização de procedimentos de referenciação lexical e pronominal que permitam a retomada de referentes textuais, como pronomes, advérbios e artigos ; além de elementos lexicais, como sinônimos, hiperônimos, hipônimos, expressões metafóricas e expressões resumitivas ) ; - emprego adequado dos pronomes relativos “ cujo(a) ” e “ onde ”, responsáveis por encaixamentos sintáticos no período composto ; - utilização adequada de sinais de pontuação que contribuam para a organização lógica da frase e do texto ; - inteligibilidade relacionada ao atendimento das exigências de estruturação textual.
A expectativa, portanto, era de que o participante recorresse a procedimentos linguístico-discursivos para organizar seu texto, permitindo o encadeamento lógico entre suas partes de forma a garantir a progressão e a coerência textuais, com estruturação que não reproduzisse hábitos da oralidade e sim atendesse à norma-padrão da Língua Portuguesa. ( 2 ) Respeito às convenções ortográficas da norma-padrão da Língua Portuguesa.
Esta competência envolve o respeito às convenções estabelecidas pela legislação em vigor e consubstanciadas no Vocabulário Ortográfico da Língua Portuguesa, editado pela Academia Brasileira de Letras e avaliam -se : - domínio das regras de acentuação gráfica ; - domínio da grafia padrão das palavras ( com ausência de abreviaturas próprias da linguagem da internet ). ( 3 ) Assim, esperava -se que o participante : grafasse corretamente as palavras ; respeitasse as regras de acentuação gráfica ; empregasse maiúsculas em início de frase, em nomes próprios de pessoas, lugares ou instituições ; grafasse as siglas com letras maiúsculas ou apenas com a primeira letra maiúscula quando formarem uma palavra ; evitasse abreviações como p/, vc, tb, pq, tá, né, usadas muitas vezes em escrita informal e na internet ; obedecesse às regras de separação de sílabas no final da linha ; atendesse aos padrões de segmentação vocabular da norma-padrão. ( 4 ) Domínio dos diferentes aspectos morfossintáticos próprios da modalidade escrita formal da norma-padrão da Língua Portuguesa.
Esta competência envolve o domínio das convenções próprias da norma-padrão da Língua Portuguesa relativas aos seguintes aspectos : concordância nominal ; concordância verbal ; regência nominal ; regência verbal ; flexão nominal ; flexão verbal ; correlação entre os tempos verbais ; colocação pronominal. ( 5 ) Assim, esperava -se que o participante : - flexionasse o verbo para estabelecer concordância de número com o sujeito da frase ; - flexionasse o artigo, o adjetivo e o pronome para concordar em número e em gênero com o substantivo a que cada um se refere ; - observasse a regência nominal e a verbal, utilizando a preposição adequada depois de um substantivo, um verbo ou um adjetivo ; - empregasse adequadamente o acento grave indicador de crase que expresse a contração entre uma preposição e um artigo, de acordo com as exigências de regência nominal ou verbal ; - obedecesse às regras de colocação pronominal ( próclise e ênclise ), distintas dos hábitos da oralidade ou da escrita informal ; - flexionasse adequadamente verbos, substantivos, adjetivos e pronomes, no que diz respeito à expressão das categorias gramaticais ; - flexionasse os verbos para expressar a correlação de modo e de tempo nas estruturas subordinadas.
Com base nesses critérios, alguns desvios ortográficos foram considerados como desvios de caráter morfossintático porque interferem na classificação dos vocábulos de base morfológica e funcional : - eliminação da marca de infinitivo ( -r- ) e substituição por acento agudo ou ausência total de marca do infinitivo, causando ambiguidade entre distintas classes de palavras ; - confusão entre as terminações “ ão ” e “ am ”, utilizadas para marcar diferentes tempos verbais ; - emprego indevido da forma verbal de"haver"\&lt ; “ há”\&gt ; no lugar da preposição ou do artigo \&lt ; “ a”\&gt ;, e vice-versa ; - emprego de hífen para separar pronome átono – tanto uso indevido quanto omissão ( ex : “ esperasse ”, em lugar de “ espera -se ” ; “ falar -mos ”, no lugar de “ falarmos ” ) ; - emprego de “ esta ” no lugar de “ está ”, provocando ambiguidade devido à confluência gráfica entre o pronome demonstrativo e a forma do verbo"estar", pela ausência da acentuação gráfica ; - emprego da conjunção aditiva"e"no lugar da forma verbal “ é ”, provocando ambiguidade devido à confluência gráfica por ausência da acentuação gráfica ; - emprego do advérbio “ mais ” no lugar da conjunção adversativa “ mas ”, reproduzindo confluência fônica da oralidade ; - registro da terceira pessoa do plural dos verbos “ ter ” e “ vir ” sem o acento circunflexo, causando desvio de concordância verbal ; - emprego inadequado do relativo “ onde ” deslocado do seu sentido locativo, reproduzindo tendência da oralidade ; - emprego do pronome relativo “ que ” sem a necessária preposição ( desvio de regência ) ; - ausência do acento grave obrigatório para marcar a crase. ( 6 ) Seleção vocabular adequada à modalidade escrita formal da Língua Portuguesa, exigida pela situação comunicativa.
Esta competência envolve : - precisão na seleção/utilização do vocabulário relacionado à temática solicitada pela questão ; - diversificação e adequação vocabular ; - coerência na combinação vocabular ; - ausência de marcas da oralidade, como termos de sentido muito genérico ( “ coisa ”, “ negócio ”, “ você ” ) e termos de registros mais informais ( como gírias, jargões, frases feitas, ditados populares, termos regionais ).
Assim, esperava -se que o participante : respeitasse a adequação vocabular, não empregando gírias ou expressões coloquiais ; apresentasse repertório vocabular diversificado e preciso ; evitasse repetição desnecessária de palavras ; utilizasse um vocabulário mais formal, como solicitado por um texto dissertativo.
Essas quatro competências foram reunidas em três itens de avaliação, atendendo aos aspectos explicitados anteriormente : I - Aspectos Ortográficos ( 30 \% ) : domínio das convenções ortográficas : grafia de vogais e consoantes, uso de maiúsculas e minúsculas, emprego do hífen e da acentuação gráfica, segmentação de palavras ; II - Aspectos Textuais ( 20 \% ) : domínio dos procedimentos de estruturação textual : organização interna dos períodos, emprego de conectores para a articulação lógica entre os períodos e entre os parágrafos, emprego de marcas de referenciação lexical e pronominal ; utilização dos sinais de pontuação que contribuem para a organização lógica da frase ; III - Aspectos Morfossintáticos e Vocabulares ( 50 \% ) : domínio das regras de caráter morfossintático estabelecidas como modelares do ponto de vista da modalidade escrita formal da norma-padrão da Língua Portuguesa : concordância nominal e verbal, regência nominal e verbal, colocação pronominal, flexão nominal e verbal, correlação entre tempos e modos verbais, ausência de marcas de oralidade.
Outra competência combinada a esta na análise das respostas é a seleção vocabular adequada à modalidade escrita formal da Língua Portuguesa, tendo em vista a intersecção entre as duas do ponto de vista das exigências do registro formal da modalidade escrita da norma-padrão. b ) Análise das respostas dadas pelos estudantes A Questão 1 propiciou a oportunidade de maior desenvolvimento do tema solicitado, gerando consequências na estruturação textual das respostas, já que os textos foram mais expositivos, mais elaborados e fluentes, inclusive, um dos textos era uma tabela.
Já a Questão 2 teve um comportamento distinto, devido ao seu caráter mais opinativo, o texto foi mais curto e o vocabulário mais limitado.
Observaram -se, nas duas questões, alguns casos de comprometimento de autoria, com textos muito próximos a paráfrases dos textos motivadores e sem marcas de repertório sociocultural produtivo.
Apresentam -se, a seguir, algumas observações sobre o desempenho dos participantes : I - Aspectos Ortográficos : o desempenho dos participantes revelou uma diferença muito grande nos dois aspectos analisados, como já foi observado nos relatórios dos anos anteriores : baixo índice de desvios da grafia padrão e grande índice de desvios de acentuação.
Em vários casos, ocorre ausência quase completa de acentuação gráfica.
Vale ressaltar que muitos participantes apresentaram uma caligrafia quase ininteligível.
Ainda nesse aspecto, ocorreu a omissão inadequada de cedilha, de til e do traço do T em algumas respostas.
Houve, inclusive, respostas com ausência total de cedilha ou til.
Os casos mais sistemáticos de eliminação do acento indicador da sílaba tônica foram : palavras proparoxítonas ( “ politica ”, “ pratica ”, “ publica ”,"policia") ; palavras paroxítonas terminadas em ditongo crescente ( “ necessario ”,"historia","propria", “ violencia ” ) ; palavras oxítonas ( “ tambem ”,"populacao", “ pais ” ).
Por outro lado, destaca -se o uso indevido do acento gráfico em palavras como “ investímento ” e “ pátriarcado ”, por exemplo.
Quanto aos desvios de grafia, observam -se os seguintes exemplos : emprego inadequado de maiúscula : “ Leis ” ; ou de minúscula :"constituição"“ Ministério publico ” ; confusão entre"ss","sc","c"e"ç": “ nessessaria", “ agreção ”, “ conciencia ”, “ comecam ”, “ cocientização ”, “ pacos ” ; eliminação de grafemas : “ mostando ”,"deveriamo"; - desvios de segmentação : “ des de", “ des igualdade ”, “ pre conceito"; troca de grafema : “ órgam ” ( por “ órgão ” ), “ escravido ” ( por “ escravidão ” ).
II - Aspectos Textuais : a estruturação textual foi a competência que apresentou o maior número de desvios, sendo responsável, em alguns textos, pela falta de clareza e pela desorganização.
No conjunto dos textos analisados, observaram -se diferentes níveis de domínio dos recursos de estruturação e organização textual.
No entanto, os participantes podem ser agrupados em três faixas descritas a seguir : - os que possuem grande domínio dos recursos de estruturação textual, com produção de textos que atendem às exigências de clareza e coerência, com progressão temática bem direcionada e utilização produtiva de operadores argumentativos ; - os que apresentam um domínio regular dos recursos de estruturação textual, com desvios na organização do texto e ausência de paragrafação, além de uso insuficiente de operadores argumentativos e processos de encaixamento sintático ; - os que apresentam baixo domínio dos recursos de estruturação textual, com muitos desvios que retomam hábitos da oralidade e que comprometem a coerência do texto : ocorrência de sequência justaposta de ideias sem encaixamentos sintáticos ; redução de estruturas subordinadas, ao lado do aumento na frequência de estruturas coordenadas e absolutas ; redução no uso de conectores para expressar relações lógicas essenciais à construção do texto ; emprego equivocado de operadores que não estabelecem relações lógicas coerentes entre ideias do texto ; fragmentação frasal que compromete a estrutura lógico-gramatical.
Uma estrutura que vem se ampliando na escrita dos participantes, refletindo hábitos da oralidade, é a duplicação de um dos elementos da frase, ou seja, a retomada anafórica de um elemento topicalizado, como por exemplo :"Elas precisam saber o verdadeiro papel DELAS na sociedade “, A desigualdade ELA vem de uma cultura machista... ” ; “ Homem e mulher, ELES são iguais ”.
Também em relação ao emprego dos sinais de pontuação, observaram -se distintos tipos de desempenho : participantes que dominam a funcionalidade da pontuação para a construção de textos fluentes e coesos, por um lado, e um grande número de participantes que apresentam precariedade em sua utilização, chegando ao ponto de elaborar parágrafos sem qualquer marca interna de pontuação para separar os períodos.
São os seguintes os tipos de desvios encontrados : - emprego de vírgula para separar o sujeito e o predicado : “ A mulher, sofre violência ” ( no lugar de “ A mulher sofre violência ” ) ;"O correto seria, ter um canal de denúncia ” ( no lugar de “ o correto seria ter um canal de denúncia ” ) ; - ocorrência de apenas uma das vírgulas para separar uma palavra, uma expressão ou uma oração encaixada : “ O homem, por sua vez não aceita que uma mulher seja sua chefa ” ( no lugar de “ O homem, por sua vez, não aceita que uma mulher seja sua chefa ” ) ; - emprego de ponto no lugar de vírgula para separar orações de um mesmo período : “ A todo custo as mulheres são atacadas na política.
Sexualmente também. ” ( no lugar de “ A todo custo as mulheres são atacadas na política, sexualmente também. ” ) ; - ausência de vírgula ou conjunção para separar elementos de uma enumeração :"... conquistas como voto equidade salarial direito de estudar e reconhecimento..."( no lugar de""... conquistas como voto, equidade salarial, direito de estudar e reconhecimento...") ; - ausência de vírgula para separar oração adjetiva explicativa ou utilização inadequada para separar oração adjetiva restritiva : “ As mulheres, que são iguais aos homens devem ser igualmente protegidas por lei."( no lugar de “ As mulheres, que são iguais aos homens, devem ser igualmente protegidas por lei.") ; - emprego indevido do sinal de dois pontos : “ Mulheres são julgadas : pelas simples roupas que usam. ” ( no lugar de “ Mulheres são julgadas pelas simples roupas que usam. ” ) ; - ausência de pronome no processo de referenciação : “ julga a mulher ou até agride ” ( no lugar de “ julga a mulher ou até a agride ” ” ) ; - emprego de vírgula em lugar do ponto-final : “ O patriarcado estrutural é formado por elementos históricos e religiosos, há pouquíssimo tempo, as mulheres passaram a votar. ” ( no lugar de “ O patriarcado estrutural é formado por elementos históricos e religiosos.
Há pouquíssimo tempo, as mulheres passaram a votar. ” ).
III - Aspectos Morfossintáticos e Vocabulares : o domínio das convenções morfossintáticas apresentou resultados que evidenciam, claramente, alguns aspectos problemáticos no desempenho dos participantes.
Assim como na análise dos desvios textuais, pode -se observar uma distinção entre três grupos de participantes : o grupo que revela domínio das convenções morfossintáticas ; o grupo que apresenta um domínio regular, com desvios pontuais ; e o grupo que revela domínio precário dessas convenções.
Pode -se destacar a regência, tanto verbal quanto nominal, como o aspecto com maior número de desvios : a ) Regência nominal e verbal Na maioria dos textos, destacam -se os desvios relacionados à noção de"regência", ou seja, à relação de subordinação que se estabelece entre as palavras por meio, ou não, de preposições, refletindo uma tendência na modalidade escrita da língua.
O desvio mais frequente foi a falta do sinal indicativo da crase, em contextos em que seria obrigatória, ou, ao contrário, um uso indevido em contextos em que o princípio não se aplica.
Os dois tipos de desvios revelam que o usuário não tem consciência de que, sob a forma do termo “ a ”, se pressupõe a presença de uma preposição “ a ”, exigida pela regência do termo anterior.
Exemplos : “ diferentes tipos de agressão a mulher ” ( no lugar de “ diferentes tipos de agressão à mulher ” ) ; “ Desenvolverem práticas econômicas e sociais as novas gerações"( no lugar de “ Desenvolverem práticas econômicas e sociais às novas gerações") ; “ Além de proteger à mulher... ” ( no lugar de “ Além de proteger a mulher... ” ).
Outro desvio relacionado à regência verbal e à nominal, encontrado frequentemente nas questões, é a ausência de preposição antes de pronome relativo, refletindo um processo generalizado na modalidade oral da língua, em situações de registro informal.
Exemplo :"outra forma é oferecer apoio psicológico quem sofreu qualquer tipo de abuso..."( no lugar de :"outra forma é oferecer apoio psicológico a quem sofreu qualquer tipo de abuso...") b ) Concordância nominal e verbal Os desvios relativos à aplicação das regras de concordância verbal e nominal correspondem a fatos que têm sido amplamente estudados na modalidade oral da Língua Portuguesa.
Quanto à concordância de número, observou -se ausência de marca ( com sujeito anteposto ou posposto ), como por exemplo a ausência de acento circunflexo na forma plural do presente do indicativo do verbo “ ter ”, considerada como um desvio na concordância verbal e não na acentuação gráfica.
Exemplo :"as mulheres ainda tem como obrigação"( no lugar de"as mulheres têm como obrigação") ; “ as leis em vigor tem que ser cumpridas ” ( no lugar de “ as leis em vigor têm que ser cumpridas ” ).
Outro desvio na aplicação da regra da concordância de número foi o aparecimento da marca de plural em verbos ou adjetivos relacionados a núcleos substantivos no singular, devido à presença de modificadores no plural, evidenciando um processo de hipercorreção.
Esse processo tem se ampliado nas situações de oralidade, caracterizando uma tendência de variação da Língua Portuguesa.
Exemplos :"Mudança nas leis que aumentem a pena do agressor"( no lugar de"Mudança nas leis que aumente a pena do agressor") ;"o aumento dos feminicídios devem ser combatidos ” ( no lugar de"o aumento dos feminicídios deve ser combatidos ” ).
Quanto à concordância de gênero, vários casos foram observados, normalmente no âmbito de sintagmas nominais longos, em que o adjetivo está afastado do substantivo, conforme apontam as pesquisas sobre a modalidade oral.
Exemplo : “ E as leis sobre as mulheres que sofrem violência ser mais rigorosos ” ( no lugar de “ E as leis sobre as mulheres que sofrem violência ser mais rigorosas ” ) ; práticas econômicas e sociais às novas gerações") ; “ Além de proteger à mulher... ” ( no lugar de “ Além de proteger a mulher... ” ).
Outro desvio relacionado à regência verbal e à nominal, encontrado frequentemente nas questões, é a ausência de preposição antes de pronome relativo, refletindo um processo generalizado na modalidade oral da língua, em situações de registro informal.
Exemplo :"outra forma é oferecer apoio psicológico quem sofreu qualquer tipo de abuso..."( no lugar de :"outra forma é oferecer apoio psicológico a quem sofreu qualquer tipo de abuso...") b ) Concordância nominal e verbal Os desvios relativos à aplicação das regras de concordância verbal e nominal correspondem a fatos que têm sido amplamente estudados na modalidade oral da Língua Portuguesa.
Quanto à concordância de número, observou -se ausência de marca ( com sujeito anteposto ou posposto ), como por exemplo a ausência de acento circunflexo na forma plural do presente do indicativo do verbo “ ter ”, considerada como um desvio na concordância verbal e não na acentuação gráfica.
Exemplo :"as mulheres ainda tem como obrigação"( no lugar de"as mulheres têm como obrigação") ; “ as leis em vigor tem que ser cumpridas ” ( no lugar de “ as leis em vigor têm que ser cumpridas ” ).
Outro desvio na aplicação da regra da concordância de número foi o aparecimento da marca de plural em verbos ou adjetivos relacionados a núcleos substantivos no singular, devido à presença de modificadores no plural, evidenciando um processo de hipercorreção.
Esse processo tem se ampliado nas situações de oralidade, caracterizando uma tendência de variação da Língua Portuguesa.
Exemplos :"Mudança nas leis que aumentem a pena do agressor"( no lugar de"Mudança nas leis que aumente a pena do agressor") ;"o aumento dos feminicídios devem ser combatidos ” ( no lugar de"o aumento dos feminicídios deve ser combatidos").
Quanto à concordância de gênero, vários casos foram observados, normalmente no âmbito de sintagmas nominais longos, em que o adjetivo está afastado do substantivo, conforme apontam as pesquisas sobre a modalidade oral.
Exemplo : “ E as leis sobre as mulheres que sofrem violência ser mais rigorosos ” ( no lugar de “ E as leis sobre as mulheres que sofrem violência ser mais rigorosas ” ) ; Registro da forma verbal “ é ” sem acento, provocando ambiguidade na leitura pela confluência com a conjunção aditiva"e".
Exemplo : “ A desvantagem entre salários, poderes e entre a violência física e moral e crescente por demais... ” ( no lugar de “ A desvantagem entre salários, poderes e entre a violência física e moral é crescente por demais ”. c ) Emprego do pronome relativo “ onde ” como relativo universal, típica marca da oralidade, em que esse pronome não é mais utilizado apenas em sua função locativa"lugar onde se realiza a ação"Exemplo :"Existe a desigualdade de gênero, onde homem ganha mais do que as mulheres"; “ há falta de informação e conhecimento onde ” ; “ O patriarcado está estruturado em nossa sociedade há séculos, onde o sexo masculino... ” d ) Falta de artigo definido antes de substantivo, comprometendo o paralelismo em enumerações Exemplo :"seja em oportunidade de emprego, salário"( no lugar de"seja em oportunidade de emprego, de salário"). e ) Desvios relacionados à flexão verbal Exemplos : emprego indevido do tempo ou do modo verbal : “ promovido ” por “ promovendo ” ; eliminação de marca de infinitivo : “ colocar ” por “ coloca ”, “ ta ” por “ estar ”, “ está ” por “ estar ”, “ esta ” por “ estar ” ; “ estás ” por “ estas ” ; utilização de construção de “ estar ” + gerúndio como marca de oralidade : “ a sociedade cada vez mais tá evoluindo ” ; segmentação indevida da desinência de número e pessoa de verbos : “ denunciar -mos ”.
Assim como nas demais competências avaliadas, há diferenças entre os participantes em relação à seleção e ao emprego do vocabulário.
Foram encontrados textos com alto padrão de domínio vocabular, chegando até ao emprego de vocabulário técnico das áreas do saber abordadas nas questões.
Por outro lado, na maioria dos textos, foram observados vários tipos de inadequação na seleção e na utilização do vocabulário, sendo que o principal deles foi a excessiva repetição de certas palavras, revelando limitação de repertório vocabular.
Além disso, observou -se a ocorrência de expressões da oralidade, revelando falta de domínio de vocabulário mais abstrato e de maior complexidade, essencial ao desenvolvimento do texto de base dissertativa.
Exemplos : seleção vocabular indevida : “ Em pleno século 21 é de emplena importância... ” ; emprego inadequado do verbo “ ter ” no lugar de “ haver ” : “ Durante a história, teve muito machismo ” ; imprecisão vocabular : “ Mas a cabeça, o corpo, há mesma igualdade para desenvolvermos igualmente as coisas ”, “ Em alguns anos atrás, as coisas eram mais complicadas nas questões de desigualdade no trabalho. ” A imprecisão vocabular, como o uso da palavra “ coisa ” em vários textos, demonstra o domínio precário da norma-padrão da Língua Portuguesa, com presença constante de traços de oralidade. 7.3.1.8 CONSIDERAÇÕES FINAIS SOBRE A CORREÇÃO DE LÍNGUA PORTUGUESA DAS RESPOSTAS ÀS QUESTÕES DISCURSIVAS DO COMPONENTE DE FORMAÇÃO GERAL O conjunto de textos analisados no Enade 2022 permite traçar um perfil bem detalhado do desempenho linguístico do estudante brasileiro de graduação. 7.3.2 COMPONENTE ESPECÍFICO A Tabela 7.17 apresenta dados sobre questões discursivas do Componente Específico.
A partir dos dados registrados na Tabela 17.7, observa -se que a “ Média ” ( 58,6 ) foi mais alta do que para as Questões Discursivas do Componente de Formação Geral.
A maior “ Média ” deste componente foi obtida pelos estudantes da região Nordeste ( 60,3 ), e a menor, pelos estudantes da região Norte ( 49,2 ).
As demais regiões obtiveram “ Média ” igual a 58,3 ( Centro-Oeste ), 60,2 ( Sudeste ) e 57,8 ( Sul ).
Quanto à variabilidade das notas, o “ Desvio-padrão ” de todo o Brasil foi 27,5.
O maior “ Desvio-padrão ” foi encontrado na região Centro-Oeste ( 29,3 ), e o menor, na região Norte ( 16,6 ).
A maior nota “ Máxima ”, 100,0, foi obtida por, pelo menos, um estudante em todas as regiões, exceto na região Norte, que obteve 71,7.
O mesmo ocorreu com a nota “ Mínima ” zero, exceto na região Norte, que obteve 28,3.
A “ Mediana ” do Brasil e da região Centro-Oeste foi igual a 63,3 ; a região Norte obteve 45,4, no Nordeste 65,0, e Sudeste foi 64,2, e na região Sul foi 62,5.
No Gráfico 7.10, representa -se a distribuição das notas nas questões discursivas do componente de específico.
Como mostra o Gráfico 7.10, o comportamento das notas das questões de formação geral apresenta um histograma com a distribuição das mesmas : uma distribuição unimodal com a moda no intervalo ( 70 ; 80 ]. 7.3.2.1 ANÁLISE DA QUESTÃO DISCURSIVA 3 DO COMPONENTE ESPECÍFICO Os resultados a Questão Discursiva 3 do componente específico são descritos na Tabela 7.18 e comentados na sequência.
Como mostrado na Tabela 7.18, na Questão 3, a “ Média ” dos estudantes de todo o Brasil foi 64,6.
A menor “ Média ” na questão, em evidência, foi obtida pelos estudantes da região Centro-Oeste ( 62,5 ), enquanto a maior “ Média ” foi obtida na região Norte ( 70,8 ).
Quanto à variabilidade das notas, o “ Desvio-padrão ” de todo o Brasil foi igual a 36,1.
O maior “ Desvio-padrão ” foi obtido na região Centro-Oeste e Sul ( 36,5 ), enquanto o menor foi obtido na região Norte ( 24,6 ).
A nota “ Máxima ”, 100,0 pontos, foi alcançada por, pelo menos, um estudante em cada região.
A “ Mediana ” do Brasil foi 75,0, para maioria das regiões ; exceto nas regiões Norte e Centro-Oeste, que foram, respectivamente, 62,5 e 50,0.
A nota “ Mínima ” zero foi a mesma em todas as regiões do Brasil, exceto na região Norte, que obteve 50,0.
No Gráfico 7.11, mostra -se a distribuição das notas na Questão Discursiva 3, do componente específico da área de Tecnologia em Design de Moda.
Como mostra o Gráfico 7.11, o comportamento das notas das Questões 3 de formação geral apresenta um histograma com a distribuição das mesmas : uma distribuição unimodal com a moda no intervalo ( 90 ; 100 ]. 7.3.2.2 COMENTÁRIOS SOBRE AS RESPOSTAS À QUESTÃO DISCURSIVA 3 A questão abordou um tema interessante para a área do Design de Moda, pois, além de desenvolver novos produtos, o designer deve ser capaz de compreender as questões estratégicas voltadas à comercialização do vestuário.
Ao contemplar uma nova forma de oferta dos produtos de moda, esperava -se que os estudantes fossem capazes de citar possíveis dificuldades encontradas pelos empreendedores no processo de divulgação, venda, entrega, troca e comunicação com o cliente, realizados de forma on-line, durante o período do isolamento social e lockdown.
Era esperado que os estudantes citassem duas dificuldades e comentassem a respeito delas, seja indicando possíveis formas de superá -las ou demonstrando os benefícios ou fracassos encontrados após a implementação de tais ações, no processo de digitalização do mercado de moda.
Entende -se que o espaço destinado à resposta foi adequado, embora grande parte das respostas não tenham ocupado todas as linhas.
A maioria dos estudantes citou, ao menos, um desafio de forma semelhante ao que foi apresentado no padrão de resposta.
Um comportamento recorrente foi que muitos participantes iniciaram a resposta contextualizando dificuldades gerais da população durante a época do isolamento social, imposto pela pandemia da Covid-19, fugindo do tema central da questão, que era citar dificuldades específicas dos empreendedores na implementação do processo digital nos negócios de moda.
Percebeu -se uma certa dificuldade na compreensão do enunciado da questão, que pedia para o estudante “ citar e comentar ” dois desafios enfrentados pelos empreendedores e, também, uma divergência no formato de apresentação das respostas.
Muitos candidatos apenas citaram os desafios em forma de tópicos, sem fazer nenhum comentário a respeito ; alguns citaram e comentaram os desafios em forma de tópicos ; outros citaram e comentaram sobre os desafios em forma de texto ; e houve, ainda, os que contextualizaram dificuldades gerais da indústria da moda, que nada tinham a ver com o processo de implementação digital do negócio e, ao fim da resposta, citaram um ou dois desafios do processo digital, sem comentar a respeito.
Com relação à linguagem utilizada, entende -se que foi suficiente para a questão.
Como a resposta não exigia o emprego de termos técnicos específicos da área, muitas respostas apresentaram linguagem discursiva informal.
De maneira geral, grande parte dos estudantes contemplou o que foi proposto pelo padrão de respostas, citando diferentes problemas enfrentados durante a implementação do processo digital nos pequenos negócios de moda.
Dentre as dificuldades citadas, a que teve maior recorrência foi a falta de conhecimento e habilidade dos empresários para lidar com as redes sociais e as ferramentas digitais, visto que muitos dos pequenos e médios negócios de moda não estavam inseridos neste mercado antes da pandemia.
Nos comentários a respeito dessa dificuldade, grande parte dos estudantes citou que houve necessidade de capacitação e que muitos empresários precisaram buscar cursos que ensinassem a utilizar as redes sociais, criar sites de e-commerce, gerenciar vendas on-line e criar conteúdo digital para atingir os clientes.
Outra dificuldade bastante apresentada nas respostas diz respeito à entrega dos produtos, pelo fato de que os pequenos negócios de moda atendiam ao público presencialmente e precisaram encontrar novas formas de fazer com que seus produtos chegassem até o cliente.
Muitos estudantes citaram a necessidade de firmar parcerias com empresas de transporte local e também em território nacional, uma vez que vender pela internet permitiu que os negócios atingissem um maior número de pessoas.
Apareceram, em menor proporção, mas de maneira bem distribuída entre as respostas, dificuldades em relação aos novos sistemas de pagamento ; necessidade de investimento em equipamentos e espaço para fotografar e divulgar os produtos ; padronização de tabela de medidas ; e questões ligadas à segurança e confiança na relação com o cliente, durante o processo de compra digital.
Alguns estudantes também pontuaram questões de comunicação com o cliente e conectividade.
Uma parte significativa dos estudantes não citou dificuldades relacionadas ao processo de digitalização dos negócios de moda, mas apontou dificuldades relacionadas ao processo de fabricação do produto, fornecimento de matéria-prima, incentivos do governo e até mesmo ao fechamento de empresas durante este período.
Embora, neste caso, a nota de avaliação tenha sido zero, por não com o processo de implementação digital do negócio e, ao fim da resposta, citaram um ou dois desafios do processo digital, sem comentar a respeito.
Com relação à linguagem utilizada, entende -se que foi suficiente para a questão.
Como a resposta não exigia o emprego de termos técnicos específicos da área, muitas respostas apresentaram linguagem discursiva informal.
De maneira geral, grande parte dos estudantes contemplou o que foi proposto pelo padrão de respostas, citando diferentes problemas enfrentados durante a implementação do processo digital nos pequenos negócios de moda.
Dentre as dificuldades citadas, a que teve maior recorrência foi a falta de conhecimento e habilidade dos empresários para lidar com as redes sociais e as ferramentas digitais, visto que muitos dos pequenos e médios negócios de moda não estavam inseridos neste mercado antes da pandemia.
Nos comentários a respeito dessa dificuldade, grande parte dos estudantes citou que houve necessidade de capacitação e que muitos empresários precisaram buscar cursos que ensinassem a utilizar as redes sociais, criar sites de e-commerce, gerenciar vendas on-line e criar conteúdo digital para atingir os clientes.
Outra dificuldade bastante apresentada nas respostas diz respeito à entrega dos produtos, pelo fato de que os pequenos negócios de moda atendiam ao público presencialmente e precisaram encontrar novas formas de fazer com que seus produtos chegassem até o cliente.
Muitos estudantes citaram a necessidade de firmar parcerias com empresas de transporte local e também em território nacional, uma vez que vender pela internet permitiu que os negócios atingissem um maior número de pessoas.
Apareceram, em menor proporção, mas de maneira bem distribuída entre as respostas, dificuldades em relação aos novos sistemas de pagamento ; necessidade de investimento em equipamentos e espaço para fotografar e divulgar os produtos ; padronização de tabela de medidas ; e questões ligadas à segurança e confiança na relação com o cliente, durante o processo de compra digital.
Alguns estudantes também pontuaram questões de comunicação com o cliente e conectividade.
Uma parte significativa dos estudantes não citou dificuldades relacionadas ao processo de digitalização dos negócios de moda, mas apontou dificuldades relacionadas ao processo de fabricação do produto, fornecimento de matéria-prima, incentivos do governo e até mesmo ao fechamento de empresas durante este período.
Embora, neste caso, a nota de avaliação tenha sido zero, por não contemplar o que era avaliado na questão, concordamos que são dificuldades que realmente existiram durante a época do lockdown e impactaram o mercado de moda.
Mesmo que as respostas tenham sido, por vezes, apenas parcialmente respondidas, apresentando, com certa frequência, dois itens com apenas um comentário, ficou evidente que os estudantes entendem a significativa importância do conhecimento para a área do Design de Moda a respeito do designer ser capaz de compreender as questões estratégicas voltadas à comercialização do vestuário.
Porém, houve dificuldade em descrever, de modo objetivo e organizado, os comentários dos itens citados, sendo, por vezes, apresentadas de forma misturada as duas justificativas, ou de forma repetida.
Identificou -se que a principal deficiência ocorreu no âmbito de interpretação de texto, no que se refere à questão que solicitava que os estudantes realizassem, durante a resposta, a seguinte ação “ cite e comente dois desafios ”.
Dessa forma, percebe -se que, em alguns casos, foi apresentada uma argumentação sobre a pergunta desnecessária ou com ausência de clareza e objetivo, não respondendo, para cada citação ou comentário, a referida questão. 7.3.2.3 ANÁLISE DA QUESTÃO DISCURSIVA 4 DO COMPONENTE ESPECÍFICO Na Tabela 7.19, constam as informações relativas à Questão 4 do conjunto de questões do componente específico.
Como se observa, a partir dos dados da Tabela 7.19, o desempenho dos estudantes nessa questão foi inferior ao da Questão 3 e superior ao da Questão 5.
A “ Média ” geral do Brasil foi 58,1, sendo a maior “ Média ” registrada na região Sudeste ( 61,1 ), e a menor, na região Norte ( 38,3 ).
Quanto à variabilidade das notas, o “ Desvio-padrão ” geral do Brasil foi 34,7, sendo o maior “ Desvio-padrão ” registrado na região Centro-Oeste ( 36,3 ), e o menor, na região Norte ( 27,9 ). contemplar o que era avaliado na questão, concordamos que são dificuldades que realmente existiram durante a época do lockdown e impactaram o mercado de moda.
Mesmo que as respostas tenham sido, por vezes, apenas parcialmente respondidas, apresentando, com certa frequência, dois itens com apenas um comentário, ficou evidente que os estudantes entendem a significativa importância do conhecimento para a área do Design de Moda a respeito do designer ser capaz de compreender as questões estratégicas voltadas à comercialização do vestuário.
Porém, houve dificuldade em descrever, de modo objetivo e organizado, os comentários dos itens citados, sendo, por vezes, apresentadas de forma misturada as duas justificativas, ou de forma repetida.
Identificou -se que a principal deficiência ocorreu no âmbito de interpretação de texto, no que se refere à questão que solicitava que os estudantes realizassem, durante a resposta, a seguinte ação “ cite e comente dois desafios ”.
Dessa forma, percebe -se que, em alguns casos, foi apresentada uma argumentação sobre a pergunta desnecessária ou com ausência de clareza e objetivo, não respondendo, para cada citação ou comentário, a referida questão. 7.3.2.3 ANÁLISE DA QUESTÃO DISCURSIVA 4 DO COMPONENTE ESPECÍFICO Na Tabela 7.19, constam as informações relativas à Questão 4 do conjunto de questões do componente específico.
Como se observa, a partir dos dados da Tabela 7.19, o desempenho dos estudantes nessa questão foi inferior ao da Questão 3 e superior ao da Questão 5.
A “ Média ” geral do Brasil foi 58,1, sendo a maior “ Média ” registrada na região Sudeste ( 61,1 ), e a menor, na região Norte ( 38,3 ).
Quanto à variabilidade das notas, o “ Desvio-padrão ” geral do Brasil foi 34,7, sendo o maior “ Desvio-padrão ” registrado na região Centro-Oeste ( 36,3 ), e o menor, na região Norte ( 27,9 ).
A nota “ Máxima ” ( 100,0 ) foi atingida por, pelo menos, um concluinte em todas as regiões.
A “ Mediana ” em todo o Brasil foi 70,0, o mesmo valor foi obtido em todas as regiões, exceto na região Norte ( 45,0 ).
Foi zero a nota “ Mínima ” de todas as regiões.
No Gráfico 7.12, está representada a distribuição de notas da Questão Discursiva 4, do Componente Específico.
Como mostra o Gráfico 7.12, o comportamento das notas das Questões 3 de formação geral apresenta um histograma com a distribuição dessas : uma distribuição unimodal com a moda no intervalo ( 0 ; 10 ]. 7.3.2.4 COMENTÁRIOS SOBRE AS RESPOSTAS À QUESTÃO DISCURSIVA 4 A questão está elaborada de forma clara e é de fácil compreensão.
Abordou um interessante tema na área do Design de Moda, que é o marketing e suas consequências negativas com a não oferta de produtos divulgados, chamando a atenção para a possível interpretação equivocada no seu entendimento de campanha e para as prováveis implicações e consequências quando uma empresa de moda não oferece o que divulga, tendo, assim, repercussão negativa, que pode causar inúmeros prejuízos, inclusive a exclusão e falência da marca.
Um tema extremamente importante para o setor da moda, que trabalha com imagem e identificação de público-alvo como um dos principais norteadores para a criação do vestuário.
Era esperado que os estudantes citassem e comentassem três consequências negativas sobre campanhas de marketing que não \textbf{ofertam} produtos que divulgam e também descrevessem duas características de produtos de vestuário que atendessem as necessidades de usuários com problemas de mobilidade.
Foi observado, no item “ a ”, que os estudantes citavam e não comentavam o quesito, conforme era pedido no enunciado, e, no item “ b ”, muitos estudantes apenas citaram, sem descreverem as características do produto.
O espaço destinado à resposta foi adequado e, em grande parte, as respostas ocuparam todas as linhas.
Os estudantes, em muitas questões, responderam igual ou de forma aproximada ao constante no padrão de resposta proposto.
No item “ a ”, foram frequentes as respostas que citavam apenas o pedido no enunciado e não comentavam as consequências negativas, conforme o pedido na questão.
No item “ b ”, ocorreu, com muita frequência, apenas a citação das características de um produto de vestuário que atendesse às necessidades de usuários com restrição de mobilidade e no enunciado o pedido era descrever.
Algumas respostas não tinham absolutamente nada a ver com o pedido no enunciado e outras repetiam a pergunta como início da argumentação, sem a sua finalização com o contexto proposto.
Uma grande parte dos estudantes respondeu conforme o padrão de respostas apresentado, muitos apenas citaram, outros citaram e comentaram sobre o tema e alguns escreveram de forma corrida contextualizando todo o enunciado, e não respondendo separadamente cada item, além de apenas citarem a reposta quando o comando era descrever.
Verificou -se uma certa dificuldade na compreensão do enunciado quando o comando pedia para “ citar e comentar ” três consequências negativas para a empresa e também descrever duas características de produtos de vestuário voltados para pessoas com restrição de mobilidade, e diversos formatos de apresentação de respostas, em tópicos ou corridas, misturadas entre os dois itens, respondidas iniciando pelo item “ b ” e também sem separação de assunto.
Vários estudantes apenas citaram as consequências em forma de tópicos, sem fazer nenhum comentário a respeito ; alguns citaram e comentaram as consequências em forma de tópicos ; outros citaram e comentaram sobre as consequências negativas em forma de texto corrido ; e, ainda, houve os que responderam sobre temas que nada tinham a ver com a pergunta em questão, bem como os que responderam ao primeiro quesito ( “ a ” ) com o tema do quesito ( “ b ” ), e o contrário também.
A linguagem utilizada pela grande maioria foi a discursiva informal, em algumas respostas, extremamente informal, mas a questão não pedia termos técnicos e específicos da área, portanto, sem prejuízos à nota.
De modo geral, grande parte dos estudantes respondeu igual ou de forma aproximada ao padrão de respostas proposto, citando e comentando as consequências negativas das campanhas de marketing e também descrevendo as características de produtos de vestuário voltados para pessoas com restrição de mobilidade.
Dentre as consequências negativas citadas, a maior ocorrência foi o cancelamento da marca, perda de clientes, e, segundo a descrição dos estudantes, a marca “ mal falada ”, imagem negativa, mas também houve ocorrências de constrangimento do cliente, perda de credibilidade, falência, processos judiciais, propaganda enganosa, dentre outras.
Com relação ao quesito “ b ”, descrição de características de produtos de vestuário voltados para pessoas com restrição de mobilidade, apareceram recorrentemente : ergonomia, tecidos leves, usabilidade, vestibilidade, tecidos maleáveis, malha, zíper, botões e aberturas laterais.
Alguns confundiram restrição de mobilidade com roupa inclusiva, embora este público entre nesse aspecto a ser incluído na moda, a pergunta relacionava -se apenas à descrição de produtos para usuários com restrição de mobilidade.
Nota -se que os estudantes têm conhecimento sobre o assunto.
Embora algumas respostas tenham sido apenas parcialmente respondidas, existe ainda uma dificuldade recorrente em expressar -se discursivamente sobre o tema e interpretar comandos, confundindo as intenções ou resumindo a maneira de explanar o conteúdo.
As respostas revelam, em parte, um problema de interpretação da pergunta, pois o comando solicitava que os estudantes citassem e comentassem as consequências negativas da realização de campanhas de marketing focadas na diversidade sem \textbf{ofertar} produtos que atendam esses usuários, item “ a ”, e, descrevessem duas características de um produto que atendessem às necessidades de usuários com restrição de mobilidade, item “ b ” ; o que ocorreu, em grande parte das respostas, foi apenas a citação, como se não houvesse o segundo comando para comentar, e, no comando para descrever, que significa “ \textbf{contar} com detalhes ”, ocorria, na sua grande maioria, apenas a citação.
Fica evidente que o estudante possui o conhecimento sobre o assunto na área do Design de Moda, embora tenha tido dificuldades em descrever objetivamente as respostas e ter misturado em muitas dela os assuntos.
A principal deficiência identificada foi na interpretação dos comandos das perguntas, em que “ citar e comentar ” eram atendidos parcialmente nas respostas, e o comando “ descrever “ era apenas citado brevemente, sem a descrição da característica de um produto que atendesse às necessidades de usuários com restrição de mobilidade.
Desta forma, alguns estudantes repetiram o enunciado e citaram algumas consequências e não comentaram sobre elas, em outros casos, davam exemplos de campanhas mal realizadas, sem citar suas consequências negativas, e, em algumas respostas, confundiram as consequências negativas de marketing com a inclusão e a restrição de movimentos. 7.3.2.5 ANÁLISE DA QUESTÃO DISCURSIVA 5 DO COMPONENTE ESPECÍFICO Na Tabela 7.20, constam as informações relativas à Questão 5 do conjunto do componente específico.
Como revelam os dados da Tabela 7.20, o desempenho dos estudantes de todo o Brasil na Questão Discursiva 5 do Componente Específico foi inferior ao das Questões Discursivas 3 e 4, deste A nota “ Média ” dos estudantes de todo o Brasil foi 53,0.
A maior “ Média ” foi registrada na região Sudeste ( 54,5 ), enquanto a menor “ Média ” foi registrada na região Norte ( 38,3 ).
Quanto à variabilidade das notas, o “ Desvio-padrão ” dos estudantes do Brasil, como um todo, foi de 31,8.
Enquanto o maior desvio foi encontrado na região Centro-Oeste ( 32,6 ), o menor foi encontrado na região Norte ( 25,6 ).
A “ Mediana ” para o Brasil e para as regiões Sudeste e Sul foi 55,0.
Para as regiões Centro-Oeste e Nordeste, a “ Mediana ” foi 50,0, e na região Norte foi 41,3.
A nota “ Máxima ” foi 100,0, obtida por, pelo menos, um estudante em cada região, exceto na região Norte que obteve 77,5.
A nota “ Mínima ” do Brasil foi zero, o mesmo valor encontrado em todas as regiões.
No Gráfico 7.13, é representada a distribuição das notas da Questão Discursiva 5 do Componente Específico.
Como mostra o Gráfico 7.13, o comportamento das notas das Questões 3 de formação geral apresenta um histograma com a distribuição dessas : uma distribuição unimodal com a moda no intervalo ( 40 ; 50 ]. 7.3.2.6 COMENTÁRIOS SOBRE AS RESPOSTAS À QUESTÃO DISCURSIVA 5 A questão abordou um tema importante para o Design de Moda, que é a estamparia digital.
Na primeira parte da pergunta, estavam relacionadas às melhorias para a empresa, sendo muitos aspectos levados em consideração no padrão de resposta, dentre elas : liberdade criativa, melhor aproveitamento do tecido, custo benefício e maiores possibilidades no uso das cores e formas.
Na segunda parte da questão, foi abordado, também, um ponto muito importante na Moda, que é a diminuição do impacto ambiental por meio da estamparia digital.
Neste ponto, foram consideradas várias possibilidades de melhoria no padrão de resposta, sendo elas voltadas para a não utilização de água, redução do gasto de energia e de descarte de materiais, tintas e tecidos no meio ambiente.
Todos os apontamentos levados em consideração nos padrões de respostas foram bem coerentes, tanto na questão A quanto na questão B, dando ao aluno possibilidades de caminhos a seguir dentro da temática proposta.
No geral, os alunos se mantiveram dentro das possibilidades de respostas apontadas no padrão, até porque o padrão de respostas era amplo, para uma melhor adequação na resposta dos alunos.
Na primeira questão, falou -se muito sobre a diminuição de custos, a facilidade e a exclusividade que a estamparia digital trazia para a empresa.
Na segunda questão, falou -se sobre a economia de água, a energia e a minimização dos produtos químicos desprezados no meio ambiente.
Com relação à linguagem utilizada pelos estudantes, as respostas estavam dentro do esperado.
Alguns estudantes utilizaram os termos técnicos de maneira muito adequada, enquanto outros usaram expressões comuns ao cotidiano ; foram detectados também muitos erros de escrita.
As respostas encontradas foram analisadas da seguinte forma, quanto ao item A : {\EmojiFont ●} Aumento do lucro – Devido à rapidez no processo de criação da estampa e à simplicidade do método, não há a necessidade de tantos recursos humanos no processo e de tantos materiais, como telas, cilindros, tintas e testes de tecidos. {\EmojiFont ●} Agilidade na produção – A agilidade do processo de impressão da estampa aumenta, consequentemente, os lucros da empresa.
Além disso, muito se falou da redução de perda de tecido por erros ou por conta dos testes, bem como da praticidade de criar no computador e a impressora transferir para o tecido, minimizando erros e aprovando a arte antes da impressão. {\EmojiFont ●} Liberdade na criação e no uso de cores – Enquanto a estampa está sendo criada, o resultado já pode ser observado, as cores podem ser alteradas e a criação pode fluir livremente, pois a quantidade de cores ou de formas não atrapalhará em nada o processo de impressão.
Muito alunos compararam muito com a serigrafia para justificar essas possiblidades de cores, formas, realismo e liberdade nas estampas.
Também foram analisadas, quanto ao item B, da seguinte forma : {\EmojiFont ●} Minimização de descartes químicos no meio ambiente – Com a estamparia digital, os impactos são minimizados devido ao não descarte no meio ambiente de pigmentos químicos, fixadores e tecidos testados. {\EmojiFont ●} Redução de lixo – Neste ponto, os alunos relacionaram a redução de lixo com a redução de materiais para a estampa, não sendo necessário desenvolver quadros e cilindros para estampas e que, posteriormente, seriam descartados no meio ambiente.
Além da diminuição do desmatamento, por não necessitar da madeira para os quadros serigráficos.
Falou -se, também, do melhor aproveitamento do tecido que diminui os resíduos têxteis e também a não necessidade de testes em tecidos, que também resultam em lixo. {\EmojiFont ●} Redução do gasto de água e de recursos naturais – Muitos alunos discorreram sobre a economia de água, visto que, neste processo, não é necessário nem lavar os tecidos nem mesmo as telas em outros métodos de estamparia.
Alguns alegaram a rapidez do processo e a necessidade de poucos maquinários resultavam também em economia de energia.
Os alunos erraram mais sobre os verbos de comando da questão ( citar e descrever 2 exemplos ), deixando algumas respostas com apenas uma citação ou descrição ou até mesmo sem se atentar que era referente ao campo da estamparia têxtil.
Porém, quanto ao conteúdo das respostas, pode -se dizer que grande parte compreendeu a técnica e a importância dela para a sustentabilidade financeira da empresa e para a sustentabilidade do Planeta.
Foi identificada uma grande deficiência na escrita dos estudantes, quando se refere à Língua Portuguesa e às pontuações.
Muitas expressões comuns na área da Moda foram escritas de forma incorreta.
Um outro aspecto relevante foi a falta de atenção dos alunos que, muitas vezes, respondiam as questões 3 ou 4 no campo da 5, sendo assim, tendo sua questão anulada. 7.3.2.7 CONSIDERAÇÕES FINAIS SOBRE AS QUESTÕES DISCURSIVAS DO COMPONENTE ESPECÍFICO No geral, o desempenho dos alunos foi muito bom, porém ficou nítido que a maioria dos estudantes não conseguiu organizar suas ideias de maneira lógica, o que ajudaria na elaboração de respostas mais precisas, aprofundadas e coerentes.
Vale ressaltar que a questão discursiva 5 foi técnica e objetiva, exigindo pouco conhecimento aprofundado dos alunos, dessa forma, as respostas também foram diretas e objetivas, com pouco aprofundamento.
É importante ampliar as discussões sobre como os estudantes de Design de Moda no Brasil estão recebendo sua formação, pois termos técnicos foram escritos de forma errada, incoerências nos conceitos e algumas justificativas inadequadas foram encontradas.
Mesmo estes sendo uma minoria, é preciso salientar que os alunos que fazem Enade estão no último semestre, prestes a se formar, e esses conceitos de Moda deveriam estar claros para eles.

% matched lemmas: contar, ofertar
\end{textsample}
